\begin{abstract}
The same compound measured against the same target in different laboratories receives different potency values, which limits how well any model can be scored against such labels. For 41 potency tasks (target $\times$ $K_i$ or $\mathrm{IC}_{50}$) from ChEMBL~37 we estimate that limit per task from 8{,}916 molecules reported in at least two source documents, and express model scores as a \emph{fraction of ceiling} $f$ (0: predicting the training mean; 1: reaching the estimated noise floor). The floors range from 0.31 to 0.92 pChEMBL units. After decontamination against MoleculeNet and TDC we benchmark ECFP-based models and frozen ChemBERTa-77M and MoLFormer-XL embeddings under four splits; the embeddings stay below ECFP on every one. The main result concerns the metric rather than the models. ECFP scores below the estimated floor on 16 of the 41 tasks under a random split, and $f$ rises with the noise it is meant to divide out ($\rho=+0.74$). The cause is document support: how much of a test molecule's own source document the model has already seen. Over 609{,}539 single-document test cases, RMSE falls from 1.07 pChEMBL with no same-document training molecule to 0.51 with a hundred or more, below the median floor of 0.54. With no support the random, scaffold and temporal splits give the same error, so their familiar ordering reflects document support rather than chemical difficulty. A constant per-document offset accounts for 57\% of the noise variance, too little to explain that gap. A document-grouped split puts every task back below the ceiling and removes the confounding ($\rho=+0.05$), while leaving the model ordering unchanged.

\textbf{Scientific Contribution.} We estimate the label-noise ceiling separately for each benchmark task from independent ChEMBL documents, instead of assuming one literature-wide noise level, and validate the estimate against held-out measurements. We show that scores at or beyond that ceiling are produced by document support, which also accounts for the apparent difficulty ordering of random, scaffold and temporal splits. We release a document-grouped split and code that computes the ceiling and the support distribution of any ChEMBL task.
\end{abstract}

\paragraph{Keywords} ChEMBL; label noise; noise ceiling; benchmarking; data splits; potency prediction; QSAR; reproducibility
